\startSection
\selectlanguage{romanian}
\sectionTitle{Imnul Național al României}

\begin{question}{Cine este autorul versurilor imnului național al României, \enquote{Deșteaptă-te, române}?}
  \choice{Vasile Alecsandri}
  \choice{Mihai Eminescu}
  \choice[correct]{Andrei Mureșanu}
  \choice{George Coșbuc}
  \choice{Ion Heliade Rădulescu}
\end{question}

\begin{question}{Cine este asociat în mod tradițional cu melodia imnului \enquote{Deșteaptă-te, române}?}
  \choice{Ciprian Porumbescu}
  \choice[correct]{Anton Pann}
  \choice{Tiberiu Brediceanu}
  \choice{George Enescu}
  \choice{Dinu Lipatti}
\end{question}

\begin{question}{În ce an a devenit \enquote{Deșteaptă-te, române} imnul național al României după schimbările politice de la sfârșitul secolului al XX-lea?}
  \choice{1918}
  \choice{1947}
  \choice{1989}
  \choice[correct]{1990}
  \choice{2001}
\end{question}

\begin{question}{Care este primul vers al imnului național românesc?}
  \choice{\enquote{Români, uniți-vă în cuget și-n simțiri}}
  \choice{\enquote{Sculați-vă din somnul cel de moarte}}
  \choice[correct]{\enquote{Deșteaptă-te, române, din somnul cel de moarte}}
  \choice{\enquote{Trăiască România liberă și unită}}
  \choice{\enquote{Ridică-te, nație mândră și curajoasă}}
\end{question}

\begin{question}{În contextul istoric al Revoluției de la 1848, ce rol simbolic are îndemnul \enquote{Deșteaptă-te}?}
  \choice{Descrie o sărbătoare agricolă.}
  \choice{Anunță începutul unei competiții sportive.}
  \choice[correct]{Cheamă la trezire civică, identitate și libertate.}
  \choice{Prezintă regulile unei ceremonii religioase.}
  \choice{Explică structura unei instituții administrative.}
\end{question}

\begin{question}{Care dintre următoarele cântece NU a fost imn național oficial al României?}
  \choice{\enquote{Trăiască Regele}}
  \choice{\enquote{Pe-al nostru steag e scris Unire}}
  \choice{\enquote{Trei culori}}
  \choice[correct]{\enquote{Hora Unirii}}
  \choice{\enquote{Deșteaptă-te, române}}
\end{question}
\shuffleSection
\renderSection